\chapter{Expectation, Covariance, Correlation \& Probability Bounds}

Summary statistics---such as the mean (center of mass), variance (dispersion), and covariance (joint variation)---condense high-dimensional joint distributions into tractable quantitative summaries. In this chapter, we master \index{Joint Expectation}\textbf{Joint Expectation}, \index{Linearity of Expectation}\textbf{Linearity of Expectation}, \index{Covariance}\textbf{Covariance}, Pearson's \index{Correlation Coefficient}\textbf{Correlation Coefficient}, and non-parametric distribution-free probability bounds via \index{Markov's Inequality}\textbf{Markov's} and \index{Chebyshev's Inequality}\textbf{Chebyshev's Inequalities}.

\vspace{0.5cm}
\begin{tcolorbox}[colback=primarylight, colframe=primarydark, title={\textbf{What You'll Learn in this Chapter}}]
\begin{itemize}
    \item Computing expectations of arbitrary bivariate functions $\E[g(X, Y)]$.
    \item Universal Linearity of Expectation (holding with or without independence).
    \item Covariance and Pearson Correlation as measures of linear association.
    \item Tail risk bounds: Markov's and Chebyshev's inequalities.
    \item Essential exam traps: $\E[g(X)] \neq g(\E[X])$ (Jensen's inequality) and Uncorrelatedness vs. Independence.
\end{itemize}
\end{tcolorbox}

\section{Expected Value and Linearity of Expectation}

\begin{definitionbox}{Expected Value of Bivariate Functions $\E[g(X, Y)]$}
For discrete random variables $X$ and $Y$ with joint PMF $p_{X,Y}(x, y)$, the expected value of any real-valued function $g(X, Y)$ is defined by the \index{Law of the Unconscious Statistician (LOTUS)}\textbf{LOTUS}:
$$\E[g(X, Y)] = \sum_{x \in \mathcal{R}_X} \sum_{y \in \mathcal{R}_Y} g(x, y) \cdot p_{X,Y}(x, y)$$
provided the absolute sum $\sum_x \sum_y |g(x, y)| p_{X,Y}(x, y) < \infty$.
\end{definitionbox}

\begin{theorembox}{Linearity of Expectation}
For \textbf{any} random variables $X_1, X_2, \dots, X_n$ (defined on the same probability space, with absolutely no requirement of independence!) and real constants $a_1, \dots, a_n, b$:
$$\E\left[ \sum_{i=1}^n a_i X_i + b \right] = \sum_{i=1}^n a_i \E[X_i] + b$$
\textbf{Special Multiplicative Case:} If $X$ and $Y$ are \textbf{independent}, then:
$$\E[X Y] = \E[X] \cdot \E[Y]$$
\end{theorembox}

\begin{videocard}{IITM BS Lecture: Expected Value \& Properties of Expectation}{https://www.youtube.com/watch?v=q6Zj54UZqYg}
Official lecture introducing expectation for joint distributions and proving linearity of expectation.
\end{videocard}

\section{Covariance and Correlation}

\begin{intuitionbox}{Visualizing Covariance as a Quadrant Balance}
Plot points $(x_i, y_i)$ with their probability masses centered at the centroid $(\mu_X, \mu_Y)$. The centroid divides the plane into four quadrants:
\begin{itemize}
    \item Quadrants I \& III: $(X - \mu_X)$ and $(Y - \mu_Y)$ share the same sign $\implies (X - \mu_X)(Y - \mu_Y) > 0$.
    \item Quadrants II \& IV: $(X - \mu_X)$ and $(Y - \mu_Y)$ have opposite signs $\implies (X - \mu_X)(Y - \mu_Y) < 0$.
\end{itemize}
$\Cov(X, Y)$ represents the probability-weighted tug-of-war between concordant and discordant quadrants.
\end{intuitionbox}

\begin{lstlisting}[language=Python, caption=Data Science: Covariance and Correlation Matrices in NumPy]
import numpy as np

# Feature vectors: Advertising Spend (k$) and Product Sales (Units)
spend = np.array([10, 20, 30, 40, 50])
sales = np.array([25, 45, 50, 70, 85])

# Bivariate Covariance Matrix (ddof=0 for population parameters)
cov_matrix = np.cov(spend, sales, ddof=0)
print("Covariance Matrix:\n", cov_matrix)
cov_xy = cov_matrix[0, 1]

# Pearson Correlation Coefficient Matrix
corr_matrix = np.corrcoef(spend, sales)
print(f"Cov(Spend, Sales): {cov_xy:.2f}")
print(f"Pearson Correlation rho: {corr_matrix[0, 1]:.4f}")
\end{lstlisting}

\begin{definitionbox}{Covariance $\Cov(X, Y)$}
The \textbf{Covariance} measures the linear co-dependence between $X$ and $Y$:
$$\Cov(X, Y) = \E[(X - \E[X])(Y - \E[Y])] = \E[XY] - \E[X]\E[Y]$$
Key Properties:
\begin{enumerate}
    \item $\Cov(X, X) = \Var(X)$.
    \item Symmetry: $\Cov(X, Y) = \Cov(Y, X)$.
    \item Bilinearity: $\Cov(aX + b, cY + d) = ac \Cov(X, Y)$.
    \item Variance of a linear combination:
    $$\Var(aX + bY) = a^2 \Var(X) + b^2 \Var(Y) + 2ab \Cov(X, Y)$$
\end{enumerate}
\end{definitionbox}

\begin{formulabox}{Pearson Correlation Coefficient $\rho_{X,Y}$}
The \index{Correlation Coefficient}\textbf{Pearson Correlation Coefficient} is the scale-invariant, dimensionless normalization of covariance:
$$\rho_{X,Y} = \frac{\Cov(X, Y)}{\sigma_X \sigma_Y} \quad \text{where } \sigma_X = \sqrt{\Var(X)}, \; \sigma_Y = \sqrt{\Var(Y)}$$
By the Cauchy-Schwarz inequality, $\rho$ is strictly bounded:
$$-1 \le \rho_{X,Y} \le 1$$
\begin{itemize}
    \item $\rho = +1 \iff Y = aX + b$ with $a > 0$ (perfect deterministic positive linear line).
    \item $\rho = -1 \iff Y = aX + b$ with $a < 0$ (perfect deterministic negative linear line).
    \item $\rho = 0 \iff X$ and $Y$ are \textbf{uncorrelated} (no linear trend).
\end{itemize}
\end{formulabox}

\begin{videocard}{IITM BS Lecture: Covariance \& Correlation Coefficient}{https://www.youtube.com/watch?v=UoZc81wNK4I}
Official lecture deriving covariance properties, variance of linear combinations, and correlation bounds.
\end{videocard}

\section{Markov and Chebyshev Inequalities}

In data science, we often work with complex black-box distributions where the true PMF/PDF is unknown, but sample estimates of the mean $\mu$ and standard deviation $\sigma$ are available. Concentration inequalities provide guaranteed, worst-case upper bounds on tail event risks.

\begin{theorembox}{Markov's Inequality}
Let $X$ be a strictly \textbf{non-negative} random variable ($X \ge 0$) with finite mean $\E[X]$. Then for any threshold $a > 0$:
$$\P(X \ge a) \le \frac{\E[X]}{a}$$
\end{theorembox}

\begin{theorembox}{Chebyshev's Inequality}
Let $X$ be any random variable (positive, negative, or mixed) with finite mean $\mu = \E[X]$ and variance $\sigma^2 = \Var(X)$. Then for any $\epsilon > 0$:
$$\P(|X - \mu| \ge \epsilon) \le \frac{\Var(X)}{\epsilon^2}$$
Equivalently, setting $\epsilon = k\sigma$ (measuring deviations in units of standard deviations $k > 0$):
$$\P(|X - \mu| \ge k\sigma) \le \frac{1}{k^2} \iff \P(|X - \mu| < k\sigma) \ge 1 - \frac{1}{k^2}$$
For instance, for any arbitrary distribution whatsoever, at least $1 - \frac{1}{2^2} = 75\%$ of all data must lie within $2\sigma$ of the mean!
\end{theorembox}

\begin{videocard}{IITM BS Lecture: Bounds in Probabilities (Markov \& Chebyshev)}{https://www.youtube.com/watch?v=NTGLrewi_uY}
Official lecture proving Markov's and Chebyshev's inequalities and applying them to tail probability bounds.
\end{videocard}

\section{Exam Gotchas \& Cognitive Traps}

\begin{trapbox}{The Nonlinear Expectation Fallacy ($\E[g(X)] \neq g(\E[X])$)}
\textbf{The Trap:} Claiming that $\E\left[\frac{1}{X}\right] = \frac{1}{\E[X]}$, or $\E[\sqrt{X}] = \sqrt{\E[X]}$, or $\E[X^2] = (\E[X])^2$.

\textbf{The Reality:} Expectation is linear, but \textbf{never passes through nonlinear functions}. 
By \index{Jensen's Inequality}\textbf{Jensen's Inequality}, for any strictly convex function $\phi$ (such as $\phi(x) = x^2$, $\phi(x) = 1/x$ for $x>0$, or $\phi(x) = e^x$):
$$\E[\phi(X)] > \phi(\E[X]) \quad \text{unless } \Var(X) = 0$$
For example, flip a fair coin where $X \in \{1, 2\}$ with $p=0.5$.
$\E[X] = 1.5 \implies \frac{1}{\E[X]} = \frac{1}{1.5} \approx 0.667$.
However: $\E\left[\frac{1}{X}\right] = \frac{1}{1}(0.5) + \frac{1}{2}(0.5) = 0.5 + 0.25 = 0.75 \neq 0.667$!
\end{trapbox}

\begin{trapbox}{Zero Covariance Does NOT Mean Independent}
\textbf{The Trap:} Concluding that because $\Cov(X, Y) = 0$ (or $\rho_{X,Y} = 0$), $X$ and $Y$ are independent.

\textbf{The Classic Counterexample:}
Let $X \in \{-1, 0, 1\}$ with uniform probability $p_X(x) = \frac{1}{3}$. Let $Y = X^2$.
\begin{itemize}
    \item Clearly, $Y$ is completely deterministic from $X$ (perfect non-linear dependence!).
    \item Compute expectations: $\E[X] = \frac{-1 + 0 + 1}{3} = 0$.
    \item $\E[XY] = \E[X^3] = \frac{(-1)^3 + 0^3 + 1^3}{3} = 0$.
    \item $\Cov(X, Y) = \E[XY] - \E[X]\E[Y] = 0 - (0)(\E[Y]) = 0$.
\end{itemize}
$\Cov(X, Y) = 0$ even though $Y$ is 100\% determined by $X$. Pearson correlation only measures \textbf{linear} relationships!
\end{trapbox}

\begin{trapbox}{Markov's Non-negativity Prerequisite}
\textbf{The Trap:} Applying Markov's inequality $\P(X \ge a) \le \frac{\E[X]}{a}$ to a random variable that can take negative values (e.g. a centered variable with mean 0).

\textbf{The Reality:} If $X$ can be negative, Markov's inequality is invalid and produces nonsensical bounds (such as negative probabilities). Markov's inequality \textbf{strictly requires $X \ge 0$}. To bound general variables, one must square the deviation $(X - \mu)^2 \ge 0$, which is precisely how Chebyshev's inequality is derived.
\end{trapbox}

\section{Comparative Synthesis Matrix}

\begin{table}[htbp]
\centering
\small
\begin{tabular}{p{3.2cm}p{3.8cm}p{3.8cm}p{3.5cm}}
\toprule
\textbf{Dimension} & \textbf{Markov's Inequality} & \textbf{Chebyshev's Inequality} & \textbf{Chernoff Bound} \\
\midrule
\textbf{Assumptions} & $X \ge 0$ (non-negative), finite $\E[X]$ & Finite mean $\mu$ and variance $\sigma^2$ & Moment Generating Function $M_X(t)$ exists \\
\addlinespace
\textbf{Target Event} & Upper tail $\P(X \ge a)$ & Two-sided deviation $\P(|X - \mu| \ge \epsilon)$ & Extreme tail $\P(X \ge a)$ \\
\addlinespace
\textbf{Decay Rate} & Hyperbolic $O(1/a)$ & Quadratic $O(1/\epsilon^2)$ & \textbf{Exponential} $O(e^{-t a})$ \\
\addlinespace
\textbf{Sharpness} & Loosest bound (needs only mean) & Moderately tight (uses variance) & Sharpest tail bounds in ML theory \\
\addlinespace
\textbf{Data Science Use} & Bounding positive resource counts (latencies, counts) & Bounding confidence intervals without normality & PAC-learning sample complexity bounds \\
\bottomrule
\end{tabular}
\caption{Comparison of Fundamental Probability Bounds and Concentration Inequalities.}
\end{table}

\newpage
\section{Practice Exercises}
\textit{Detailed step-by-step solutions are available in the Appendix.}

\subsection*{Part A: Analytical \& Computational Problems}

\begin{enumerate}
\item \textbf{Linearity of Expectation:} Roll 10 fair 6-sided dice. Let $X_i$ be the score on die $i$, and $S = \sum_{i=1}^{10} X_i$ be the total sum. Find $\E[S]$.

\item \textbf{Covariance Calculation from Joint PMF:} The joint PMF of $X \in \{0, 1\}$ and $Y \in \{1, 2\}$ is:
    $$p_{X,Y}(0,1) = 0.2, \quad p_{X,Y}(0,2) = 0.3, \quad p_{X,Y}(1,1) = 0.1, \quad p_{X,Y}(1,2) = 0.4$$
    Compute $\E[X]$, $\E[Y]$, $\E[XY]$, and the covariance $\Cov(X, Y)$.

\item \textbf{Correlation Coefficient:} For the variables $X$ and $Y$ in Exercise 2:
    \begin{enumerate}
        \item Compute $\Var(X)$ and $\Var(Y)$.
        \item Compute the Pearson correlation coefficient $\rho_{X,Y}$.
    \end{enumerate}

\item \textbf{Markov's Inequality Tail Bound:} The average daily website visits for a blog is $\E[X] = 500$. Using Markov's inequality, find an upper bound for the probability that on a given day the blog receives at least 2,000 visits ($\P(X \ge 2000)$).

\item \textbf{Data Science Application (Algorithmic Trading Risk Bound):} An automated trading strategy has expected daily return $\mu = \$100$ with standard deviation $\sigma = \$20$. Use Chebyshev's inequality to find a lower bound for the probability that daily return falls between $\$60$ and $\$140$.
\end{enumerate}

\subsection*{Part B: Conceptual Diagnostic Traps (Exam Gotchas)}

\begin{enumerate}[resume]
\item \textbf{Diagnostic Trap (Nonlinear Expectation \& Jensen's Inequality):}
Let $X$ be a non-degenerate positive random variable ($\Var(X) > 0$ and $X > 0$). Which of the following relationships is mathematically guaranteed?
\begin{enumerate}
    \item $\E\left[\frac{1}{X}\right] = \frac{1}{\E[X]}$
    \item $\E\left[\frac{1}{X}\right] > \frac{1}{\E[X]}$
    \item $\E\left[\frac{1}{X}\right] < \frac{1}{\E[X]}$
    \item $\E[\ln X] > \ln(\E[X])$
\end{enumerate}

\item \textbf{Diagnostic Trap (Uncorrelatedness and Independence):}
Suppose two continuous features satisfy $\Cov(X, Y) = 0$. Which of the following conclusions is strictly valid?
\begin{enumerate}
    \item $X$ and $Y$ are statistically independent.
    \item $p_{X,Y}(x, y) = p_X(x) p_Y(y)$ for all $(x, y)$.
    \item $\Var(X - Y) = \Var(X) + \Var(Y)$.
    \item $\E[X^2 Y^2] = \E[X^2] \E[Y^2]$.
\end{enumerate}
\end{enumerate}
